% Irodov Part V §5.7, originals 269--280, printed p.230.
% Source: user-supplied Arihant PDF p.239; work functions from Appendix 22
% (PDF p.378), zinc=3.74 eV, copper=4.47 eV, cesium=1.89 eV, lithium=2.39 eV.

\thermalopticsproblem{269}{230}
一颗质量为 $M$、半径为 $R$ 的恒星表面发出
角频率为 $\omega_0$ 的光子，求远处观测者接收光子的
相对引力频移 $\Delta\omega/\omega_0$。
\exsolution
在恒星外部近似采用史瓦西静态时空。
静态观察者测得的局部频率满足
$\omega(r)\sqrt{1-2GM/(rc^2)}=\mathrm{const}$。
远处的频率因而是
\[
 \omega_\infty=\omega_0\sqrt{1-\frac{2GM}{Rc^2}}.
\]
故
\[
 \boxed{\frac{\Delta\omega}{\omega_0}
 =\sqrt{1-\frac{2GM}{Rc^2}}-1
 \simeq-\frac{GM}{Rc^2}},
\]
后一式用于 $2GM/(Rc^2)\ll1$。
\excheck 光子向外克服引力势，远处频率\emph{降低}。
题目只给 $M,R$，不可臆造具体数值。

\thermalopticsproblem{270}{230}
X 射线管电压升高到原来的
$\eta=1.5$ 倍，连续谱短波限移动
$\Delta\lambda=26\,\mathrm{pm}$。
求原电压 $U_1$。
\exsolution
连续谱短波限由单电子能量全部转为光子给出
$\lambda_{\min}=hc/(eU)$。
原、后波长分别为
$\lambda_1=hc/(eU_1)$、
$\lambda_2=hc/(e\eta U_1)$，
故
\[
 \Delta\lambda=\lambda_1-\lambda_2
 =\frac{hc}{eU_1}\left(1-\frac1\eta\right).
\]
整理得到
\[
 \boxed{U_1=\frac{hc}{e\Delta\lambda}
 \left(1-\frac1\eta\right)
 \simeq15.9\,\mathrm{kV}}.
\]
\excheck 升高电压时短波限向\emph{更短}波长移动；
$\Delta\lambda$ 是减少量的正值。

\thermalopticsproblem{271}{230}
窄 X 射线束照射 NaCl 单晶，
可观察到镜面布拉格反射的最小掠射角为
$\alpha=4.1^\circ$，
晶面距 $d=0.28\,\mathrm{nm}$。
求射线管的加速电压。
\exsolution
满足一级布拉格反射的最短可用波长为
$\lambda_{\min}=2d\sin\alpha$；
它也必须等于连续谱截止波长
$hc/(eU)$。
因此
\[
 \boxed{U=\frac{hc}{2ed\sin\alpha}
 \simeq31.0\,\mathrm{kV}}.
\]
\excheck 所用 $\alpha$ 是相对晶面的\emph{掠射角}，
不是相对法线的入射角；截止处取布拉格一级 $m=1$。

\thermalopticsproblem{272}{230}
X 射线管内电子到达靶前速度为 $v=0.85c$。
求轫致辐射连续谱的最短波长。
\exsolution
电子为相对论粒子，其动能为
$K=(\gamma-1)m_ec^2$，
$\gamma=(1-v^2/c^2)^{-1/2}$。
最短波长来自全部动能转为单光子：
\[
 \boxed{\lambda_{\min}=
 \frac{hc}{(\gamma-1)m_ec^2}}.
\]
用 $\gamma\simeq1.8983$、
$m_ec^2\simeq511\,\mathrm{keV}$，
\[
 \boxed{\lambda_{\min}\simeq2.70\,\mathrm{pm}
 \approx2.8\,\mathrm{pm}}.
\]
\excheck $v=0.85c$ 时不能用非相对论
$m_ev^2/2$ 代替电子动能。

\thermalopticsproblem{273}{230}
求锌的光电效应截止波长，
以及波长 $250\,\mathrm{nm}$ 光照下
锌表面光电子的最大速度。
\exsolution
原书附录22给锌的逸出功
$A_{\mathrm{Zn}}=3.74\,\mathrm{eV}$。
截止处 $hc/\lambda_{\mathrm{th}}=A$，所以
\[
 \boxed{\lambda_{\mathrm{th}}=
 \frac{hc}{A_{\mathrm{Zn}}}\simeq332\,\mathrm{nm}}.
\]
入射光子能量为
$hc/(250\,\mathrm{nm})\simeq4.96\,\mathrm{eV}$，
故最大电子动能
$K_{\max}\simeq1.22\,\mathrm{eV}$，满足
$K_{\max}=m_ev_{\max}^2/2$，所以
\[
 \boxed{v_{\max}=
 \sqrt{\frac{2}{m_e}
 \left(\frac{hc}{250\,\mathrm{nm}}-A_{\mathrm{Zn}}\right)}
 \simeq6.55\times10^5\,\mathrm{m/s}}.
\]
\excheck 两个光子能量均按 eV 一致换算；
锌的截止波长大于入射波长，确实发生光电发射。

\thermalopticsproblem{274}{230}
同一金属分别以
$\lambda_1=0.35\,\mu\mathrm m$、
$\lambda_2=0.54\,\mu\mathrm m$ 的光照射，
两次产生的光电子最大速度之比为
$\eta=v_1/v_2=2.0$。
求该金属逸出功 $A$。
\exsolution
爱因斯坦光电方程给
$\tfrac12m_ev_i^2=hc/\lambda_i-A$。
由动能比为速度比平方，
\[
 \frac{hc/\lambda_1-A}
 {hc/\lambda_2-A}=\eta^2.
\]
解得
\[
 \boxed{A=
 \frac{\eta^2hc/\lambda_2-hc/\lambda_1}
 {\eta^2-1}
 \simeq1.88\,\mathrm{eV}\approx1.9\,\mathrm{eV}}.
\]
\excheck 不能把速度比 $\eta$ 当作能量比，
应使用 $\eta^2=4$。

\thermalopticsproblem{275}{230}
一个远离其他导体的铜球受到
波长 $\lambda=140\,\mathrm{nm}$ 的辐射。
在无漏电、弱场近似下，球最多能被充到多少正电势？
\exsolution
铜的逸出功按原书附录22为
$A_{\mathrm{Cu}}=4.47\,\mathrm{eV}$。
光电子最大动能
$K_{\max}=hc/\lambda-A_{\mathrm{Cu}}$。
随着球失去电子，其正电势增高，
出射电子还须克服 $e\varphi$；
停止发射条件为 $e\varphi_{\max}=K_{\max}$。
故
\[
 \boxed{\varphi_{\max}=
 \frac1e\left(\frac{hc}{\lambda}-A_{\mathrm{Cu}}\right)
 \simeq4.39\,\mathrm V\approx4.4\,\mathrm V}.
\]
\excheck 球在光电发射后带\emph{正电}；
若 $hc/\lambda<A_{\mathrm{Cu}}$，则不能仅靠此光致正充电。

\thermalopticsproblem{276}{230}
一束照射锂表面的调幅电磁辐射，
电场为
$E(t)=a(1+\cos\omega t)\cos\omega_0t$，
其中 $\omega=6.0\times10^{14}\,\mathrm{s^{-1}}$，
$\omega_0=3.60\times10^{15}\,\mathrm{s^{-1}}$。
求逸出光电子的最大动能。
\exsolution
先将调制信号展开：
\[
 E(t)=a\cos\omega_0t+
 \frac a2\cos(\omega_0+\omega)t+
 \frac a2\cos(\omega_0-\omega)t.
\]
三个谱线的角频率为
$\omega_0,\omega_0\pm\omega$，
最大光子能量对应高频边带 $\omega_0+\omega$。
锂逸出功为原书附录22的
$A_{\mathrm{Li}}=2.39\,\mathrm{eV}$，因此
\[
 \boxed{K_{\max}=
 \hbar(\omega_0+\omega)-A_{\mathrm{Li}}
 \simeq0.375\,\mathrm{eV}\approx0.38\,\mathrm{eV}}.
\]
\excheck 最大动能由\emph{最高频的频谱分量}决定，
不是由电场振幅大小决定。

\thermalopticsproblem{277}{230}
波长 $\lambda=0.30\,\mu\mathrm m$ 的光照射
工作于饱和状态的光电管，
其光谱灵敏度为
$J=4.8\,\mathrm{mA/W}$。
求每个入射光子平均产生的光电子数（量子产额）。
\exsolution
设入射功率为 $P$，
每秒入射光子数 $P/(hc/\lambda)$。
设量子产额为 $w$，饱和光电流
$I=e wP/(hc/\lambda)$。
因此
\[
 \boxed{w=\frac{Jhc}{e\lambda}}.
\]
代 $J=4.8\times10^{-3}\,\mathrm{A/W}$、
$\lambda=3.0\times10^{-7}\,\mathrm m$，得
\[
 \boxed{w\simeq1.98\times10^{-2}\approx2.0\%}.
\]
\excheck 饱和电流时每个逸出电子均被收集，
这是用测得的电流推求量子产额的条件。

\thermalopticsproblem{278}{230}
真空光电管两极板分别由铯与铜制成，
外电路短接。用波长
$\lambda=0.22\,\mu\mathrm m$ 的光照射铯极，
求到达铜极的光电子最大速度。
\exsolution
铯与铜的逸出功在原书附录22中分别为
$A_{\mathrm{Cs}}=1.89\,\mathrm{eV}$、
$A_{\mathrm{Cu}}=4.47\,\mathrm{eV}$。
电子自铯表面出射最大动能为
$K_{\mathrm{Cs}}=hc/\lambda-A_{\mathrm{Cs}}$。
由于两极板\emph{电路短接}，
接触电势使从铯到铜的电子能量变化
为 $A_{\mathrm{Cs}}-A_{\mathrm{Cu}}$，
所以铜极前的最大动能为
\[
 K_{\mathrm{Cu}}=
 K_{\mathrm{Cs}}+A_{\mathrm{Cs}}-A_{\mathrm{Cu}}
 =\frac{hc}{\lambda}-A_{\mathrm{Cu}}.
\]
因此
\[
 \boxed{v_{\max}=\sqrt{\frac{2}{m_e}
 \left(\frac{hc}{\lambda}-A_{\mathrm{Cu}}\right)}
 \simeq6.4\times10^5\,\mathrm{m/s}}.
\]
\excheck 虽由铯产生光电子，
短路两电极间的接触势差不能忽略；
最终能量由\emph{铜}的逸出功控制。

\thermalopticsproblem{279}{230}
真空光电管锌阴极受到波长
$\lambda=262\,\mathrm{nm}$ 的辐射，
当外加减速电压 $U_{\mathrm{ext}}=1.5\,\mathrm V$
时光电流刚被阻断。
求光电管的接触电势差大小及其相对外电压的极性。
\exsolution
锌逸出功为 $A_{\mathrm{Zn}}=3.74\,\mathrm{eV}$，
由光电方程
\[
 K_{\max}=\frac{hc}{262\,\mathrm{nm}}-A_{\mathrm{Zn}}
 \simeq0.993\,\mathrm{eV}.
\]
电子实际遇到的总反向电势差的大小应为
$U_{\mathrm{stop}}=K_{\max}/e\simeq0.993\,\mathrm V$。
外加反向电压为 $1.5\,\mathrm V$，故内部接触势
必须抵消其一部分：
\[
 \boxed{|U_{\mathrm c}|=
 U_{\mathrm{ext}}-U_{\mathrm{stop}}
 \simeq0.51\,\mathrm V\approx0.5\,\mathrm V}.
\]
\excheck 接触电势的极性与\emph{外加减速电压相反}，
否则光电流会在低于 $1.5\,\mathrm V$ 的外电压下被阻断。

\thermalopticsproblem{280}{230}
利用光速 $c$、粒子质量 $m$ 和普朗克常数 $h$，
构造具有长度量纲的物理量，并说明其名称。
\exsolution
若长度尺度为 $L=h^\alpha m^\beta c^\gamma$，
依 SI 量纲
$[h]=ML^2T^{-1}$、
$[m]=M$、$[c]=LT^{-1}$，
要求 $M^0L^1T^0$，得到
$\alpha=1,\beta=-1,\gamma=-1$。
故
\[
 \boxed{\lambda_C=\frac{h}{mc}},
\]
它是质量为 $m$ 的粒子的\emph{康普顿波长}。
\excheck 此处使用普通普朗克常数 $h$；
约化康普顿波长则是 $\hbar/(mc)=\lambda_C/(2\pi)$。
